\section{Discussion}\label{sec:discussion}

\pseudopara{Broader impacts.} Research into LLM jailbreaks is useful for understanding the weaknesses of large language models and identifying vulnerabilities that vendors can patch. This kind of research is often described as ``red-teaming,'' widely acknowledged as an important step in deploying foundation models. The latest draft of the EU AI Act suggests red-teaming as one useful way of validating the safety of foundation models~\cite{eu2023aiActAmendment102}, and the White House's Voluntary AI Commitments include a commitment to red-teaming from seven major model vendors~\citep{whitehouse2023voluntarycommitments, whitehouse2023factsheet}. We hope that our benchmark will help researchers in this area better evaluate the misuse potential of new jailbreak techniques and thus focus resources on the most important vulnerabilities.

\pseudopara{Limitations.}\label{sec:limitations} We note three limitations with our current work. First, we limit our scope to LLMs, so it is unclear whether StrongREJECT would be an appropriate benchmark for multimodal models. 

Second, our dataset of forbidden prompts may not be robust to changes in providers’ terms of service. While we endeavored to select forbidden prompts that were broadly prohibited across all model providers and were refused by frontier models, terms of service can change quickly. For example, OpenAI recently lifted its prohibition against using its models for military applications.

Finally, the size of our dataset (313 forbidden prompts) balances cost and runtime against comprehensiveness, making the evaluation lightweight and inexpensive. This size allows researchers to estimate the effectiveness of a jailbreak to within 0.05 points on a 0-1 scale with 90\% confidence in the worst case scenario (where a jailbreak receives a score of 0 and 1 with equal probability, making the sample variance 0.25). Specifically, this sample size enables estimation to within $1.64 * \sqrt{0.25 / 313} = 0.046$ points with 90\% confidence.

However, we acknowledge that the relatively small size of our dataset is a limitation of our work. While useful for jailbreak algorithm developers who need to evaluate their attacks frequently during development, StrongREJECT's 313-prompt dataset is not large enough to be a stand-alone, definitive measure of model robustness. Some of the best NLP benchmarks have thousands of examples; for instance, MMLU contains 15,908 questions. Therefore, we emphasize that for model developers seeking a more comprehensive view of model robustness, StrongREJECT should be complemented with additional analyses and larger-scale evaluations.

\pseudopara{Conclusion.} Most jailbreak papers claim the jailbreaks they introduce are highly effective. However, we show that these claims are often significantly exaggerated, making it impossible for researchers to compare jailbreaks and hindering jailbreak research. We suggest that this problem arises because of the lack of a standardized high-quality benchmark. In absence of such a benchmark, researchers have evaluated their jailbreaks using datasets of forbidden prompts that are often vague, ill-posed, unanswerable, or not actually forbidden, and evaluation methods that systematically exaggerate jailbreak effectiveness. To address these issues, this paper introduces a high-quality jailbreak benchmark--StrongREJECT--which includes a dataset of high-quality forbidden prompts and an automated evaluator that achieves state-of-the-art agreement with human judges of jailbreak effectiveness.

Notably, we find that many existing evaluation methods significantly overstate jailbreak effectiveness compared to human judgments and our StrongREJECT evaluator. We discover that this discrepancy is due to a surprising phenomenon: jailbreaks bypassing a victim model’s safety fine-tuning tend to reduce its ability to provide an attacker with specific, harmful information. These results further underscore the need for jailbreak researchers to use evaluators that, like StrongREJECT, consider both a model’s willingness to respond to forbidden prompts as well as its capability of doing so. High-quality evaluators like StrongREJECT allow researchers to distinguish between ineffective jailbreaks and those, like PAP and PAIR, which present the most serious threats.